\question{2.1}{
    Aan $20$ personen die allemaal $40$ jaar oud zijn, is gevraagd hoeveel jaar zij volledig dagonderwijs hebben genoten sinds hun eerste levensjaar.
    De resultaten waren als volgt:
    \begin{center}
        \begin{tabular}{cccccccccc}
            \multicolumn{10}{l}{\textit{Aantal jaren dagonderwijs van 20 personen}} \\
            \toprule
                $10$ & $16$ & $18$ & $16$ & $13$ & $21$ & $18$ & $20$ & $12$ & $15$ \\
                $12$ & $9$ & $14$ & $10$ & $15$ & $12$ & $19$ & $12$ & $13$ & $10$ \\
            \bottomrule
        \end{tabular}
    \end{center}
}

\begin{enumerate}[label=(\alph*)]
    \item Bereken voor deze groep van $20$ personen het rekenkundig gemiddelde van het aantal jaren op school.
    \answer{
        Voor de berekening van het rekenkundig gemiddelde van $20$ waarden geldt de volgende formule:
        \begin{align*}
            \samplemeanx = \frac{x_1 + x_2 + \ldots + x_n}{n} = \frac{x_1 + x_2 + \ldots + x_{20}}{20}
        \end{align*}

        De $20$ uitkomsten leveren gesommeerd:
        \begin{align*}
            \sum_i x_i = x_1 + x_2 + \ldots + x_{20} = 10 + 16 + \ldots + 10 = 285
        \end{align*}
        Voor het rekenkundig gemiddelde vinden we dan:
        \begin{align*}
            \samplemeanx = \frac{x_1 + x_2 + \ldots + x_{20}}{20} = \frac{285}{20} = 14.25.
        \end{align*}
    }

    \item Bereken de mediaan.
    \answer{
        Om de mediaan te kunnen bepalen, zetten we de getallen op volgorde.
        Dit levert:
        \begin{center}
            \begin{tabular}{cccccccccc}
                \multicolumn{10}{l}{\textit{Aantal jaren dagonderwijs van 20 personen}} \\
                \toprule
                    $9$ & $10$ & $10$ & $10$ & $12$ & $12$ & $12$ & $12$ & $13$ & $13$ \\
                    $14$ & $15$ & $15$ & $16$ & $16$ & $18$ & $18$ & $19$ & $20$ & $21$ \\
                \bottomrule
            \end{tabular}
        \end{center}
        Omdat we hier een even aantal uitkomsten hebben, geldt als mediaan het gemiddelde van de twee middelste uitkomsten, oftewel:
        \begin{align*}
            \text{mediaan } = \frac{13 + 14}{2} = 13.5.
        \end{align*}
    }

    \item Bepaal de modus.
    \answer{
        In het geval van afzonderlijke waarnemingen zoals hier, is de modus gedefinieerd als de uitkomst met de hoogste frequentie.
        Inspectie van de uitkomsten levert ons dat de uitkomst $12$ het meeste voorkomt. De modus is dus $12$, deze komt namelijk vier keer voor, meer dan welke andere uitkomst.
    }
    
\end{enumerate}      

